\documentclass[11pt]{article}
\usepackage[margin=1in]{geometry}
\usepackage{enumitem}
% =============================================================================
% Preambles/header.tex — shared LaTeX/Beamer preamble for this project
%
% Use from a Beamer lecture by prepending:
%     \input{header}
% and compiling with TEXINPUTS=../Preambles:$TEXINPUTS (the /compile-latex
% skill does this automatically).
%
% PALETTE CONTRACT. Color names defined here MUST match the names in
% Quarto/theme-template.scss so Beamer and Quarto renderings use the same
% palette. The scripts/check-palette-sync.sh script enforces this.
%
% To customize: edit the HEX values below AND the matching $variable in
% Quarto/theme-template.scss, then run scripts/check-palette-sync.sh.
% =============================================================================

% -----------------------------------------------------------------------------
% Engine detection + encoding
%
% Our /compile-latex skill uses XeLaTeX, where inputenc is unnecessary and
% can produce encoding warnings. Only load inputenc under pdfTeX.
% -----------------------------------------------------------------------------
\usepackage{iftex}
\ifPDFTeX
  \usepackage[utf8]{inputenc}
\fi

\usepackage{amsmath, amssymb, amsthm}
\usepackage{booktabs}
\usepackage{graphicx}

% -----------------------------------------------------------------------------
% PALETTE — mirrors Quarto/theme-template.scss variable names.
% Load xcolor and define colors BEFORE anything that references them
% (notably hyperref, hypersetup, and the Beamer color assignments).
% -----------------------------------------------------------------------------
\usepackage{xcolor}

\definecolor{primary-blue}{HTML}{012169}      % headings, accents
\definecolor{primary-gold}{HTML}{B9975B}      % emphasis, borders
\definecolor{highlight-yellow}{HTML}{F2A900}  % markers, alerts
\definecolor{light-bg}{HTML}{E8EDF5}          % title / section backgrounds
\definecolor{jet}{HTML}{1A1A1A}               % body text

% Semantic colors (used in diagrams and annotations)
\definecolor{positive}{HTML}{15803D}          % good / observed / identified
\definecolor{negative}{HTML}{B91C1C}          % bad / problematic / bias
\definecolor{neutral}{HTML}{525252}           % reference / context

% Highlight accents (match .hi-* classes in the SCSS)
\definecolor{hi-slate}{HTML}{314F4F}
\definecolor{hi-green}{HTML}{2E8B57}
\definecolor{hi-red}{HTML}{C41E3A}

% Semantic aliases so skill templates and snippets can write
% \textcolor{accent}{...} without hard-coding "primary-blue".
\colorlet{accent}{primary-blue}
\colorlet{accent2}{primary-gold}

% -----------------------------------------------------------------------------
% Hyperref — loaded AFTER xcolor + palette so link colors resolve.
% -----------------------------------------------------------------------------
\usepackage{hyperref}
\hypersetup{colorlinks=true, linkcolor=primary-blue, urlcolor=primary-blue,
            citecolor=primary-blue, pdfborder={0 0 0}}

% -----------------------------------------------------------------------------
% Course code — set once per deck (after \input{header}, before \begin{document})
% via \coursecode{CS401}. Shown in the footer on every slide so a deck's course
% is identifiable without opening the title page. Empty by default.
% -----------------------------------------------------------------------------
\makeatletter
\newcommand{\coursecode}[1]{\gdef\@coursecodetext{#1}}
\gdef\@coursecodetext{}
\makeatother

% -----------------------------------------------------------------------------
% Beamer theme assignments (only apply under Beamer).
%
% We detect Beamer via \@ifclassloaded{beamer}, which is the documented,
% non-internal way. This must be wrapped in \makeatletter/\makeatother
% because \@ifclassloaded contains an @-catcode character.
% -----------------------------------------------------------------------------
\makeatletter
\@ifclassloaded{beamer}{%
  \usefonttheme{professionalfonts}
  \setbeamercolor{structure}{fg=primary-blue}
  \setbeamercolor{frametitle}{fg=primary-blue}
  \setbeamercolor{title}{fg=primary-blue}
  \setbeamercolor{subtitle}{fg=primary-gold}
  \setbeamercolor{author}{fg=primary-blue}
  \setbeamercolor{institute}{fg=neutral}
  \setbeamercolor{date}{fg=primary-gold}
  \setbeamercolor{itemize item}{fg=primary-blue}
  \setbeamercolor{itemize subitem}{fg=primary-gold}
  \setbeamercolor{alerted text}{fg=primary-gold}
  \setbeamercolor{block title}{fg=white, bg=primary-blue}
  \setbeamercolor{block body}{bg=light-bg}
  % Semantic block variants: block=definition/concept (blue), exampleblock=
  % worked example (green/positive), alertblock=key takeaway or caution (gold).
  % Keep to <=2 colored boxes per slide across ALL three variants (INV-7).
  \setbeamercolor{block title example}{fg=white, bg=positive}
  \setbeamercolor{block body example}{bg=light-bg}
  \setbeamercolor{block title alerted}{fg=white, bg=primary-gold}
  \setbeamercolor{block body alerted}{bg=light-bg}

  % Minimal footer: course code (left) + page X / N (right), no navigation clutter
  \setbeamertemplate{navigation symbols}{}
  \setbeamertemplate{footline}{%
    \color{neutral}\footnotesize\hspace{0.7em}\@coursecodetext\hfill
    \insertframenumber/\inserttotalframenumber\hspace{0.7em}\vspace{0.3em}}
}{}
\makeatother

% -----------------------------------------------------------------------------
% TikZ setup — libraries the snippet gallery uses
% -----------------------------------------------------------------------------
\usepackage{tikz}
\usetikzlibrary{arrows.meta, positioning, calc, decorations.pathreplacing,
                fit, shapes.geometric, backgrounds}

% Shared TikZ styles — snippets and hand-written diagrams can pick these up
% via `\begin{tikzpicture}[every node/.style={...}]` or by naming specific
% styles (e.g., `\node[dag-node] (X) at (0,0) {X};`).
\tikzset{
  % Node styles for causal DAGs and similar
  dag-node/.style = {
    draw=primary-blue, fill=light-bg, rounded corners=2pt,
    minimum width=1.2cm, minimum height=0.9cm, align=center, font=\small
  },
  decision-node/.style = {
    draw=primary-gold, fill=white, diamond, aspect=1.8,
    minimum width=1.8cm, minimum height=1.2cm, align=center, font=\small,
    inner sep=1pt
  },
  % Edge styles — observed vs counterfactual
  observed-edge/.style = {-{Latex[length=2.5mm]}, thick, draw=jet},
  counterfactual-edge/.style = {-{Latex[length=2.5mm]}, thick, draw=neutral, dashed},
  confound-edge/.style = {-{Latex[length=2.5mm]}, thick, draw=negative, dashed},
  % Data-point styles for plots
  observed-dot/.style = {fill=primary-blue, draw=primary-blue, circle, inner sep=1.2pt},
  counterfactual-dot/.style = {fill=white, draw=neutral, circle, inner sep=1.2pt}
}

% -----------------------------------------------------------------------------
% Convenience macros
% -----------------------------------------------------------------------------
\newcommand{\muted}[1]{\textcolor{neutral}{#1}}
\newcommand{\key}[1]{\textcolor{primary-gold}{\textbf{#1}}}
\newcommand{\good}[1]{\textcolor{positive}{#1}}
\newcommand{\bad}[1]{\textcolor{negative}{#1}}

% Standout frame macro: full-bleed dark background for section transitions.
% Beamer-only (uses \begin{frame}/\setbeamercolor/\usebeamerfont) -- do not
% call from an article-class file (e.g. Notes/<CODE>/*-notes.tex). \muted,
% \key, \good, \bad, and \coursecode above are plain xcolor/gdef and are
% safe to use from either class; verified when Notes/article-class support
% was added.
% Expand: \transitionslide{Your transition title}
\newcommand{\transitionslide}[1]{%
  \begingroup
  \setbeamercolor{background canvas}{bg=primary-blue}
  \begin{frame}[plain]
    \centering
    \vfill
    {\usebeamerfont{title}\color{white}\Large #1\par}
    \vfill
  \end{frame}
  \endgroup
}

% Section divider frame (Beamer-only), an alternative to \transitionslide with a
% darker two-line style: near-black (jet) background, a small white label line
% above a larger gold title, and a thin gold rule along the bottom edge.
% Expand: \sectiondivider{Section 2}{Pointers}
\newcommand{\sectiondivider}[2]{%
  \begingroup
  \setbeamercolor{background canvas}{bg=jet}
  \begin{frame}[plain]
    \vspace*{3.0cm}
    {\color{white}\ttfamily\large #1\par}
    \vspace{0.5cm}
    {\color{primary-gold}\LARGE #2\par}
    \begin{tikzpicture}[remember picture, overlay]
      \draw[primary-gold, line width=2.5pt]
        ($(current page.south west)+(0,0.1)$) -- ($(current page.south east)+(0,0.1)$);
    \end{tikzpicture}
  \end{frame}
  \endgroup
}

% -----------------------------------------------------------------------------
% Channel watermark — "Concepts That Click" (YouTube).
%
% Article-class files (CompetitiveExam Q/A sets, Notes, Assignments) get a
% light diagonal draft watermark on every page plus a centered footer naming
% the channel. Beamer decks get a subtle TikZ background watermark instead
% (the footline already carries the course code). Centralized here so every
% MCQ PDF is branded without per-file edits.
% -----------------------------------------------------------------------------
\makeatletter
\@ifclassloaded{article}{%
  \usepackage{draftwatermark}
  \SetWatermarkText{Concepts That Click}
  \SetWatermarkScale{1}
  \SetWatermarkFontSize{2cm}
  \SetWatermarkLightness{0.86}
  \SetWatermarkAngle{45}
  \usepackage{fancyhdr}
  \pagestyle{fancy}
  \fancyhf{}
  \fancyfoot[C]{\footnotesize\textcolor{neutral}{Concepts That Click~|~YouTube: \href{https://www.youtube.com/@concepts.thatclick}{@concepts.thatclick}}}
  \renewcommand{\headrulewidth}{0pt}
  \renewcommand{\footrulewidth}{0pt}
  \fancypagestyle{plain}{%
    \fancyhf{}%
    \fancyfoot[C]{\footnotesize\textcolor{neutral}{Concepts That Click~|~YouTube: \href{https://www.youtube.com/@concepts.thatclick}{@concepts.thatclick}}}%
    \renewcommand{\headrulewidth}{0pt}%
    \renewcommand{\footrulewidth}{0pt}%
  }
}{}
\@ifclassloaded{beamer}{%
  \setbeamertemplate{background}{%
    \begin{tikzpicture}[remember picture, overlay]
      \node[rotate=45, opacity=0.055, align=center]
        at (current page.center)
        {\Huge\textcolor{neutral}{Concepts That Click}};
    \end{tikzpicture}%
  }
}{}
\makeatother 
\coursecode{JKSSB}
\renewcommand{\thesection}{08.\arabic{section}}
\newtheorem{example}{Example}[section]
\newenvironment{definitionbox}[1]{%
  \par\noindent\textbf{Definition (#1).}\ \itshape
}{\par\vspace{0.3em}}
\title{RAM, ROM \& Cache --- Detailed Lecture Notes}
\author{JKSSB Computer Awareness}
\date{\today}

\begin{document}
\maketitle

\section{Memory, volatility, and why the distinction matters}
\textbf{Main memory} (also called primary memory) holds the data and
instructions that the CPU uses directly. Secondary storage such as a hard
disk or an SSD is not accessed by the CPU directly; its contents are brought
into main memory first. This lesson concentrates on the primary-memory
components a computer-awareness paper actually asks about: RAM, ROM, and
cache.

The single most tested distinction in this lesson is \textbf{volatility}.

\begin{definitionbox}{Volatile memory}
Memory whose stored contents are lost when the power supply is removed.
\end{definitionbox}

\begin{definitionbox}{Non-volatile memory}
Memory whose stored contents remain available after the power supply is
removed.
\end{definitionbox}

RAM is volatile; ROM is non-volatile. A practice item may phrase this as a
NOT/incorrect-statement question, so learn both directions: ``RAM loses its
contents on power-off,'' and ``ROM retains its contents without power.''

\section{RAM: volatile main memory}
\textbf{RAM} stands for \textbf{Random Access Memory}. It is the main memory
in which running programs and their active data are held. The term ``random
access'' means that any storage location can be reached in roughly the same
amount of time; it is not a statement about volatility. RAM is
\textbf{volatile}: when the office desktop (EX-01) is switched off, whatever
was in RAM is lost unless it was saved to secondary storage.

\begin{example}
A stem states that RAM keeps its data intact after the computer is switched
off. This is false. RAM is volatile, so its contents are lost on power-off;
it is ROM (and flash) that retains data without power. If the stem instead
asks which memory holds the running programs, select RAM.
\end{example}

\section{DRAM versus SRAM}
RAM comes in two main types, and the exam routinely contrasts them. Both are
RAM and both are volatile; the difference lies in speed, refresh, cost, and
use.

\begin{center}
\begin{tabular}{p{0.22\textwidth}p{0.33\textwidth}p{0.33\textwidth}}
\toprule
\textbf{Feature} & \textbf{DRAM} & \textbf{SRAM}\\
\midrule
Full form & Dynamic Random Access Memory & Static Random Access Memory\\
Refresh & Needs periodic refresh & No refresh needed\\
Speed & Slower & Faster\\
Density and cost & Denser and cheaper per bit & Less dense and costlier\\
Typical use & Main memory (the computer's RAM) & Cache memory\\
\bottomrule
\end{tabular}
\end{center}

\textbf{DRAM} stores each bit in a tiny capacitor that leaks charge, so it
must be \textbf{refreshed} periodically. That refresh overhead is the price
of its high density and low cost, which is why DRAM is the natural choice
for large main memory. \textbf{SRAM} uses a different cell design that holds
its state without refresh, making it faster but more expensive and less
dense; it is therefore used for the small, fast \textbf{cache}.

\begin{example}
An item asks which RAM type is used to build cache memory. SRAM is faster
and needs no refresh, so it is used for cache; DRAM is denser and cheaper
and serves as main memory. If the stem asks which type needs periodic
refreshing, the answer is DRAM.
\end{example}

\section{ROM and the ROM family}
\textbf{ROM} stands for \textbf{Read-Only Memory}. It is \textbf{non-volatile}
and is \textbf{read-only in normal operation}: the processor reads its
contents, and changing them requires a special programming process rather
than ordinary memory writes. ROM commonly stores \textbf{firmware}, such as
the start-up instructions a computer runs before an operating system loads.

\begin{definitionbox}{ROM}
Read-Only Memory: non-volatile primary memory that is normally read but not
written during ordinary operation.
\end{definitionbox}

Several ROM variants appear in exams. They differ mainly in \emph{when} they
are programmed and \emph{how} they are erased.

\begin{center}
\begin{tabular}{p{0.16\textwidth}p{0.74\textwidth}}
\toprule
\textbf{Type} & \textbf{Key property}\\
\midrule
PROM & Programmable Read-Only Memory: can be programmed \textbf{once} by the user; it cannot be erased and reprogrammed.\\
EPROM & Erasable Programmable Read-Only Memory: erased by exposing the chip to \textbf{ultraviolet (UV)} light.\\
EEPROM & Electrically Erasable Programmable Read-Only Memory: erased and rewritten \textbf{electrically}, without removing the chip.\\
Flash & An \textbf{EEPROM-derived} non-volatile memory erased in \textbf{blocks}; widely used in USB drives and SSDs.\\
\bottomrule
\end{tabular}
\end{center}

\begin{example}
An item asks which memory can be erased by ultraviolet light. PROM is
one-time programmable and cannot be erased; EEPROM is erased electrically;
flash is erased in blocks. The UV-erasable type is EPROM. If the stem
instead asks which type is programmable only once, select PROM.
\end{example}

\section{Cache memory and its levels}
\textbf{Cache} is a small, very fast memory placed between the CPU and main
memory. It holds frequently or recently used data and instructions so the
CPU does not have to wait as often for slower main memory. Cache is
\textbf{faster than main memory} --- a fact the paper has tested directly
(PYQ-03). It is temporary working memory, \textbf{not} a backup and not a
substitute for storage.

Cache is organised in levels.

\begin{definitionbox}{L1 cache}
The level-1 cache: the smallest, fastest cache, located closest to the CPU
core.
\end{definitionbox}

\begin{definitionbox}{L2 cache}
The level-2 cache: larger and slightly slower than L1.
\end{definitionbox}

\begin{definitionbox}{L3 cache}
The level-3 cache: the largest and slowest of the three levels, typically
shared across the cores of the processor.
\end{definitionbox}

As the level number rises, size grows and speed falls, and distance from the
core increases: L1 $<$ L2 $<$ L3 in size, with the speed order reversed.
Even L3, however, is faster than main memory.

\section{Cache hits, misses, and hit ratio}
Performance depends on how often the CPU finds what it needs already in the
cache.

\begin{definitionbox}{Cache hit}
A memory access in which the required data is found in the cache.
\end{definitionbox}

\begin{definitionbox}{Cache miss}
A memory access in which the required data is not found in the cache, so it
must be fetched from main memory.
\end{definitionbox}

The \textbf{hit ratio} measures how often accesses succeed in the cache:
\[
\text{hit ratio} = \frac{\text{hits}}{\text{total accesses}}.
\]

\begin{example}
If the cache receives 100 memory accesses and 90 of them are found in the
cache, then the hit ratio is $90/100 = 0.90$. The remaining 10 accesses are
misses and must go to main memory. A higher hit ratio means the CPU waits
less, so the system performs better.
\end{example}

\section{Exam retrieval and common confusions}
Use the following distinctions to answer both recall and NOT/incorrect
items reliably:
\begin{itemize}
  \item RAM is \textbf{volatile}; ROM is \textbf{non-volatile} (TRAP-01).
  \item ROM is \textbf{read-only in normal operation}; writing needs a
    special process (TRAP-02).
  \item \textbf{Cache is faster} than main memory, not slower (PYQ-03).
  \item \textbf{SRAM} is used for cache; \textbf{DRAM} is used for main
    memory --- and DRAM is the type that needs refresh.
  \item Cache is \textbf{not} a backup; it is temporary working memory.
  \item ROM family: PROM (once), EPROM (UV), EEPROM (electrical), flash
    (block erase).
  \item Hit ratio $=$ hits $\div$ total accesses.
\end{itemize}

\subsection*{Exercises}
\begin{enumerate}[label=\arabic*.]
  \item Explain why RAM is described as volatile and ROM as non-volatile.
  \item Distinguish DRAM from SRAM in terms of speed, refresh, and typical
    use.
  \item Expand PROM, EPROM, EEPROM, and flash, and state how each is erased
    (or whether it can be).
  \item Order L1, L2, and L3 cache by size and by speed, and state which
    level is typically shared across cores.
  \item A cache handles 200 memory accesses and 150 are hits. Compute the
    hit ratio.
\end{enumerate}

\subsection*{Solutions}
\begingroup\small
\begin{enumerate}[label=\arabic*.]
  \item RAM is volatile because it loses its contents when power is
    removed; ROM is non-volatile because it retains its contents without
    power.
  \item DRAM needs periodic refresh, is slower but denser and cheaper, and
    is used for main memory. SRAM needs no refresh, is faster but costlier
    and less dense, and is used for cache.
  \item PROM: Programmable Read-Only Memory, programmable once and not
    erasable. EPROM: Erasable Programmable Read-Only Memory, erased by
    ultraviolet light. EEPROM: Electrically Erasable Programmable Read-Only
    Memory, erased electrically. Flash: EEPROM-derived, erased in blocks.
  \item By size: L1 $<$ L2 $<$ L3, so L3 is the largest. By speed: L1 is
    fastest and L3 the slowest of the three. L3 is typically shared across
    cores.
  \item Hit ratio $= 150/200 = 0.75$.
\end{enumerate}
\endgroup
\end{document}
